\section{Implementation and Experimental Methodology}
\label{sec:implementation-methodology}

This section describes the concrete implementation in \system and the evaluation protocol used to ensure reproducible end-to-end comparison.

\subsection{Implementation in \system}

\textbf{Execution model.}
\system is implemented in modern C++ on a shared-memory multicore runtime with eager processing. Each arriving record triggers immediate state update and opposite-side probe, followed by exact temporal/similarity verification before emission.

\textbf{Strategy composition.}
Join behavior is instantiated from four coupled dimensions: join algorithm, partition strategy, window-state type, and index strategy. The same runtime hosts baseline methods (BruteForce, IVF, HNSW, HDR-Tree, LSH, ClusteredJoin, S3J-style) and the VSJoin-oriented path.

\textbf{VSJoin routing and indexing path.}
The VSJoin implementation uses LSH-based partitioning, optional boundary-aware multicast, and logical partition virtualization. A record is first mapped to one or multiple physical partitions and then expanded to logical partitions via a stable identity hash. Logical partitions are resolved to workers through an assignment table with atomic publish semantics. Candidate retrieval probes both local mutable indexes and global read-oriented indexes, then deduplicates by record identity.

\textbf{Background maintenance.}
The runtime executes a periodic control loop for VSJoin maintenance: (i) evaluate skew from sampled load signals and apply bounded logical remapping when imbalance exceeds threshold, and (ii) rebuild global indexes from state snapshots using event-time-valid windows. The rebuild path skips semantically empty ticks to avoid churn on empty snapshots.

\textbf{Concurrency semantics.}
Index operations are mediated by a concurrency manager. Routing-map publication is read-optimized: readers do not block on updates, and writers batch changes before atomic swap. Rebalancing updates affect future routing decisions and do not require synchronous full-state migration in the foreground path.

\subsection{Workloads and Configuration Space}

We evaluate both synthetic and dataset-backed stream settings with sliding windows. Each run is defined by:
\begin{itemize}
    \item stream rates and inter-arrival characteristics,
    \item window parameters ($W$ and step size),
    \item similarity threshold $\theta$,
    \item method-specific index and routing parameters,
    \item runtime parallelism $P$.
\end{itemize}

For VSJoin-specific analysis, we sweep the parameters most tied to recent implementation changes:
\begin{itemize}
    \item multicast fanout (or equivalent boundary fanout controls),
    \item logical virtualization factor $V$,
    \item rebalance imbalance threshold $\tau$,
    \item per-round migration cap $M$,
    \item global rebuild interval.
\end{itemize}

\subsection{Metrics and Reporting}

We report three metric groups:
\begin{itemize}
    \item \textbf{Quality}: recall and precision against exact or strong-reference join outputs,
    \item \textbf{Efficiency}: throughput and end-to-end latency (including tail behavior),
    \item \textbf{Robustness}: skew sensitivity and stability under evolving stream distributions.
\end{itemize}

Latency is measured along the full operator path rather than isolated index queries. Throughput and latency are jointly interpreted to avoid over-optimizing one dimension at the expense of quality.

\subsection{Experimental Protocol}

To ensure fair comparison across methods, all algorithms run in the same runtime and share identical upstream/downstream execution semantics. Unless otherwise stated:
\begin{itemize}
    \item hardware and compiler settings are fixed across methods,
    \item each configuration is repeated across multiple runs,
    \item reported values are aggregated using robust summaries (e.g., median and spread),
    \item ablations change one design factor at a time.
\end{itemize}

We structure experiments into three blocks:
\begin{enumerate}
    \item \textbf{End-to-end method comparison}: VSJoin path versus baselines under matched window/query semantics.
    \item \textbf{Design ablation}: isolate the effect of two-tier indexing, boundary-aware fanout, and logical remapping.
    \item \textbf{Stress and sensitivity}: evaluate robustness under skew, rate bursts, and parameter sweeps for $\tau$, $M$, $V$, and rebuild interval.
\end{enumerate}

This protocol is designed to separate algorithmic gains from implementation artifacts and to make the impact of each mechanism measurable in isolation.
